\begin{thebibliography}{99}
\addcontentsline{toc}{section}{References}

\bibitem{repo:main}
ProveIt Contributors,
\emph{Polylogarithms and their Arithmetic Bridges}, canonical manuscript,
main file and preface. Repository snapshot \pin.
\url{https://github.com/VladimirReshetnikov/ProveIt/tree/73f18ef013d5b74eafc883efb2c15210642df12a/Analysis/Polylogarithms/docs/manuscript}.

\bibitem{repo:dhr}
ProveIt Contributors,
\emph{Determinantal Harmonic Resonance: Gauss--Vandermonde finite parts,
Stieltjes jets, and depth-two polylogarithmic closure}, preserved report,
\texttt{article\_standalone.tex}, at the same snapshot.
\url{https://github.com/VladimirReshetnikov/ProveIt/tree/73f18ef013d5b74eafc883efb2c15210642df12a/Analysis/Polylogarithms/docs/reports/stieltjes-harmonic-resonance/continuations/determinantal-harmonic-resonance}.
The repeated-pole, harmonic-rank, and rational-grid questions in its
further-research section are the direct targets of this continuation.

\bibitem{repo:incoming}
ProveIt Contributors,
\emph{Incoming research packages}, directory catalogue and selected
status/readme material, inspected 11 October 2026 at \pin.
\url{https://github.com/VladimirReshetnikov/ProveIt/tree/73f18ef013d5b74eafc883efb2c15210642df12a/docs/incoming}.
The precise limits of interior inspection are stated in Section~\ref{sec:audit}.

\bibitem{repo:integration}
ProveIt Contributors,
\emph{Integration chapter},
canonical manuscript \texttt{chapters/07-integration.tex}, opening 280
source lines inspected at the same snapshot.

\bibitem{repo:discovery}
ProveIt Contributors,
\emph{Discovery and proof-status chapter}, canonical manuscript
\texttt{chapters/10-discovery.tex}, at the same snapshot.

\bibitem{chen-singer}
S. Chen and M. F. Singer,
\emph{Residues and Telescopers for Rational Functions},
arXiv:1203.4200 (2012).
\url{https://arxiv.org/abs/1203.4200}.

\bibitem{arreche-sitaula}
C. E. Arreche and H. P. Sitaula,
\emph{A computational approach to rational summability and its applications
via discrete residues}, arXiv:2503.15636 (2025).
\url{https://arxiv.org/abs/2503.15636}.
Cited for the broader discrete-residue and rational-summability framework,
not as a source of the determinant identities proved here.

\bibitem{dlmf:hurwitz}
NIST Digital Library of Mathematical Functions,
\S25.11, \emph{Hurwitz Zeta Function}; in particular the defining series,
shift and parameter-derivative formulas, polygamma values, and Lerch's
identity at spectral order zero.
\url{https://dlmf.nist.gov/25.11}.

\bibitem{dlmf:gamma}
NIST Digital Library of Mathematical Functions,
Chapter 5, \emph{Gamma Function}, especially \S5.2 and \S5.15.
\url{https://dlmf.nist.gov/5.2}; \url{https://dlmf.nist.gov/5.15}.

\bibitem{dlmf:gauss}
NIST Digital Library of Mathematical Functions,
\S15.4, equation 15.4.20, Gauss hypergeometric summation.
\url{https://dlmf.nist.gov/15.4.E20}.

\bibitem{dlmf:polylog}
NIST Digital Library of Mathematical Functions,
\S25.12, \emph{Polylogarithms}, and \S25.14, \emph{Lerch's Transcendent}.
\url{https://dlmf.nist.gov/25.12}; \url{https://dlmf.nist.gov/25.14}.
The multiple-polylogarithm convention is defined explicitly in
\eqref{eq:MPL}.

\end{thebibliography}
